% titel: Blatt 10, A6: pH schwacher Säuren (Chloressigsäuren)
% skript: Mortimer Kap. 17 (schwache Säuren, pKs)

% Teil I: Werkzeuge, kein Lösungsweg.

Bei einer schwachen Säure stellt sich das Gleichgewicht
$\mathrm{HA}+\mathrm{H_2O}\rightleftharpoons\mathrm{A^-}+\mathrm{H_3O^+}$
ein. Für verdünnte Lösung gilt die Näherung.

\begin{werkzeugkasten}[pH schwacher Säure]
Gegeben $c_0$ und $pK_s$:
\[
  pH \approx \frac{1}{2}\left(pK_s - \lg c_0\right)
  \qquad(\text{falls } K_s \ll c_0).
\]
Exakt (falls nötig): $[H^+]^2 + K_s[H^+]-K_s c_0 = 0$. Elektronenziehende
Substituenten (hier Cl am $\alpha$-C) erhöhen die Acidität deutlich:
jede Chlorgruppe senkt $pK_s$ (stärkere Säure). Die Zahl der Cl
entspricht der Aciditäts-Reihenfolge Gras acetate bis Trichloracetat.
\textit{Skript:} Mortimer Kap. 17 (schwache Säure, pKs).
\end{werkzeugkasten}

\begin{achtung}[Darauf achten]
\begin{itemize}
  \item Je kleiner der $pK_s$, desto stärker die Säure $\Rightarrow$ desto
        kleiner der pH bei gleichem $c_0$.
  \item Die Näherung $\frac12(pK_s-\lg c_0)$ gilt, solange die Säure
        schwach dissoziiert.
  \item Lewis-Strukturen: Carboxylgruppe $-\mathrm{COOH}$, die Cl-Atome am
        $\alpha$-C (CH$_2$, CHCl, CCl$_2$, CCl$_3$).
  \item Trichloressigsäure ($pK_s\approx0.52$) ist fast stark sauer; die
        Näherung bleibt meist noch ok.
\end{itemize}
\end{achtung}
